\section{Discussion}
\label{sec:discussion}

\paragraph{What the representation is.}
We set out to learn one representation that serves material, force and contact across months, and found that the three need different references. Appearance carries the day signature and needs the loaded state removed; marker geometry is nearly untouched by rendering drift and needs only the unloaded grid, refreshed at a gel swap; contact needs the unloaded state too. A robot hand can maintain both states from its own stream: the loaded state from its recent grasps, the unloaded state from the bare gel between them. We state this as a finding rather than as a single network, because every attempt to learn force from appearance---including with the force signal in the objective---failed across days.

\paragraph{Limitations.}
Our own recordings come from one sensor unit; the signature and the calibration transfer to a second unit in public data, the label-free swap does not transfer to coarse shape there. Each session holds one material, so material identity can only be tested across campaigns, and the five materials recorded on the new gel are also those with a single training session, so the gel effect on material recognition cannot be isolated; it is shown through the state fields, the force scale and contact instead. Absolute accuracies are modest for a 20-way zero-shot problem; with 10\pct\ of target-day labels, supervised models still win, so the advantage of state removal is zero-shot. Lubricant is re-applied in step with use, so slow drift with use cannot be separated from lubricant residue.

\paragraph{Outlook.}
The day signature exists in every vision-based tactile dataset recorded in sittings, which is all of them; it is invisible to within-collection evaluation and to background subtraction. Reporting a later-session hold-out, and estimating the sensor state from contacts rather than from blank frames, are cheap changes that we hope become standard.
